\documentclass[11pt]{article}
\usepackage[T1]{fontenc}
\usepackage[margin=1in]{geometry}
\usepackage{amsmath,amssymb,amsthm}
\usepackage{enumitem}
\usepackage{hyperref}

\newtheorem{theorem}{Theorem}
\newtheorem{lemma}{Lemma}
\newtheorem{remark}{Remark}

\newcommand{\E}{\mathbb E}
\newcommand{\N}{\mathbb N}
\newcommand{\dd}{\,d}
\newcommand{\sfT}{\mathsf T}
\newcommand{\sfP}{\mathsf P}
\newcommand{\sfF}{\mathsf F}

\title{A Lacunary Martingale-Model Endpoint for the Subordinated Poisson Square Function}
\author{fa\_banach\_001, agent\_lane\_18}
\date{2026-07-18}

\begin{document}
\maketitle

\section*{Source Question}

The source paper is Zhendong Xu and Hao Zhang,
\emph{From the Littlewood-Paley-Stein Inequality to the
Burkholder-Gundy Inequality}, arXiv:2111.05164.  Its introduction records that
Xu had proved
\[
  \mathsf L_p^{\sfP}\lesssim \max\{p,\sqrt{p'}\}
\]
for the reverse Littlewood-Paley-Stein inequality of the Poisson semigroup
subordinated to a symmetric diffusion semigroup, and that the corresponding
problem for \(\mathsf L_p^{\sfP}\) remains open at the endpoint \(p\to1\).
Equivalently, Xu's Problem 8.4 in arXiv:2105.12175 asks, in particular,
whether for a symmetric submarkovian or Markovian semigroup there is a
semigroup-dependent constant \(C\) such that
\[
  \|f-\sfF f\|_{L^1(\Omega)}
  \le C\|G^{\sfP}_2(f)\|_{L^1(\Omega)} .
\]

\section*{Result Type}

This packet is a partial result.  It proves the endpoint estimate for a
martingale-generated model class, with sufficiently lacunary generator
spectrum.  It does not settle the arbitrary symmetric submarkovian case.

\section*{Definitions}

Let \((\Omega,\mathcal A,\mu)\) be a probability space equipped with an
increasing filtration \((\mathcal A_n)_{n\ge1}\).  Let \(E_n\) denote conditional
expectation onto \(\mathcal A_n\), set \(D_n=E_n-E_{n-1}\) for \(n\ge2\), and
write \(F=E_1\).  For a martingale terminal value \(f\), define
\[
  S(f-Ff)=\left(\sum_{n\ge2}|D_n f|^2\right)^{1/2}.
\]

Let \(0=\lambda_1<\lambda_2<\lambda_3<\cdots\), \(\lambda_n\to\infty\), and
define on martingale differences
\[
  \sfT_t f = Ff+\sum_{n\ge2} e^{-\lambda_n t}D_n f,\qquad t>0.
\]
This is the Xu-Zhang martingale construction from arXiv:2111.05164, written in
generator form: choosing \(a_{n-1}=1-e^{-\lambda_n}\) gives their
\(\sfT_t=T^t\).  In particular \((\sfT_t)_{t>0}\) is a symmetric diffusion
semigroup.  Its subordinated Poisson semigroup is
\[
  \sfP_t f=Ff+\sum_{n\ge2}e^{-t\sqrt{\lambda_n}}D_n f.
\]
The Poisson \(g\)-function is
\[
  G^{\sfP}(f)
  =\left(\int_0^\infty\left|t\frac{\partial}{\partial t}\sfP_t f\right|^2
       \frac{\dd t}{t}\right)^{1/2}.
\]

\section*{Proof Intuition}

On each martingale difference \(D_n f\), the Poisson derivative contributes the
scalar kernel \(t\sqrt{\lambda_n}e^{-t\sqrt{\lambda_n}}\).  If the numbers
\(\sqrt{\lambda_n}\) are sufficiently separated, these kernels form a Riesz
sequence in \(L^2((0,\infty),dt/t)\).  Therefore the Poisson square function is
pointwise equivalent to the usual martingale square function.  Davis'
endpoint square-function inequality then gives the desired \(L^1\) reverse
estimate.

\begin{lemma}[Lacunary Poisson kernels]
\label{lem:gram}
Let \(\alpha_{n+1}/\alpha_n\ge \rho>9\), and define
\[
  h_n(t)=2t\alpha_n e^{-t\alpha_n},\qquad t>0,
\]
as elements of \(H=L^2((0,\infty),dt/t)\).  Put
\(\eta=8/(\rho-1)<1\).  Then for every finitely supported scalar sequence
\((c_n)\),
\[
  (1-\eta)\sum_n |c_n|^2
  \le \left\|\sum_n c_n h_n\right\|_H^2
  \le (1+\eta)\sum_n |c_n|^2 .
\]
\end{lemma}

\begin{proof}
The diagonal entries are normalized:
\[
  \|h_n\|_H^2=4\alpha_n^2\int_0^\infty t e^{-2\alpha_n t}\,\dd t=1.
\]
For \(m<n\),
\[
  \langle h_m,h_n\rangle_H
  =4\alpha_m\alpha_n\int_0^\infty t e^{-(\alpha_m+\alpha_n)t}\,\dd t
  =\frac{4\alpha_m\alpha_n}{(\alpha_m+\alpha_n)^2}.
\]
Since \(\alpha_n/\alpha_m\ge \rho^{n-m}\), this is at most
\(4\rho^{-(n-m)}\).  Hence the sum of absolute off-diagonal entries in every
row of the Gram matrix is at most
\[
  2\sum_{k\ge1}4\rho^{-k}=\frac{8}{\rho-1}=\eta .
\]
The conclusion follows from the elementary Gershgorin/Schur bound for this
Hermitian Gram matrix.
\end{proof}

\begin{theorem}[Endpoint for lacunary martingale Poisson models]
\label{thm:endpoint}
Assume \(\sqrt{\lambda_{n+1}/\lambda_n}\ge\rho>9\) for all \(n\ge2\).
Then, pointwise for all finite martingales,
\[
  \frac{\sqrt{1-\eta}}{2}\,S(f-Ff)
  \le G^{\sfP}(f)
  \le \frac{\sqrt{1+\eta}}{2}\,S(f-Ff),
  \qquad \eta=\frac{8}{\rho-1}.
\]
Consequently there is a constant \(C_\rho\), independent of the filtration and
of \(f\), such that
\[
  \|f-Ff\|_{L^1(\Omega)} \le C_\rho\|G^{\sfP}(f)\|_{L^1(\Omega)} .
\]
In particular, in this model class the reverse Poisson constants remain
bounded at the endpoint \(p\downarrow1\).
\end{theorem}

\begin{proof}
For a fixed \(\omega\in\Omega\), put \(c_n=D_n f(\omega)\) and
\(\alpha_n=\sqrt{\lambda_n}\).  Since
\[
  t\frac{\partial}{\partial t}\sfP_t f(\omega)
  =-\sum_{n\ge2}t\alpha_n e^{-t\alpha_n}D_n f(\omega),
\]
we have
\[
  G^{\sfP}(f)(\omega)
  =\frac12\left\|\sum_{n\ge2}c_n h_n\right\|_{L^2(dt/t)} .
\]
Lemma~\ref{lem:gram} gives the asserted pointwise comparison with
\((\sum |D_n f(\omega)|^2)^{1/2}\).

Davis' endpoint martingale square-function inequality gives a universal
constant \(C_D\) such that
\[
  \|f-Ff\|_1\le C_D\|S(f-Ff)\|_1
\]
for every scalar martingale.  Combining this with the lower pointwise bound
for \(G^{\sfP}\) yields
\[
  \|f-Ff\|_1
  \le \frac{2C_D}{\sqrt{1-\eta}}\|G^{\sfP}(f)\|_1 .
\]
The finite-martingale proof extends to the usual \(L^1\) domain by truncating
the martingale and applying monotone convergence to the square functions.
The bounded \(p\downarrow1\) statement follows from the same
Davis-Burkholder-Gundy inequalities, whose reverse square-function constants
are bounded on \(1\le p\le2\).
\end{proof}

\section*{Verification Notes}

The only nontrivial semigroup input is the Xu-Zhang construction: increasing
\(\lambda_n\) corresponds to an increasing sequence
\(a_{n-1}=1-e^{-\lambda_n}\), so their Theorem 1.1 gives a symmetric diffusion
semigroup with eigenvalue \(e^{-\lambda_n t}\) on \(D_nL^p\).  The Poisson
subordination then replaces \(\lambda_n\) by \(\sqrt{\lambda_n}\), exactly as
used above.

The lacunarity threshold \(\rho>9\) is not intended to be sharp.  It is chosen
only to make the Riesz-sequence estimate a one-line Gershgorin argument.
A sharper Toeplitz or total-positivity analysis may cover much denser spectra,
including the numerical regime of Xu-Zhang's displayed \(16^n\) heat rates, but
that strengthening is not needed for this packet.

\section*{Limitations}

This packet does not prove Xu's full Problem 8.4.  The proof uses the
martingale difference spectral representation and a lacunarity hypothesis on
the generator.  Arbitrary symmetric diffusion semigroups need not have such a
discrete martingale-difference model, and the general semigroup-BMO dual route
remains open in this work.

\section*{Novelty and Literature Check}

Local searches of the run registry, solution index, attempt index, and proof
gap index for arXiv:2111.05164, Littlewood-Paley-Stein, Poisson reverse
constants, martingale cotype, and symmetric diffusion semigroups found no prior
packet for this target.  Local source checks included arXiv:2111.05164,
arXiv:2105.12175, and arXiv:2105.12795.  Web searches on 2026-07-18 for
phrases including ``\(L_p^P\)'', ``\(p\to1\)'', ``reverse
Littlewood-Paley-Stein'', ``Poisson semigroup'', and ``semigroup BMO'' found
Xu's explicit Problem 8.4 and classical Euclidean positive results, but no
recorded theorem proving the lacunary martingale-model endpoint above.

\section*{References}

\begin{enumerate}[label={[\arabic*]}]
\item Z. Xu and H. Zhang, \emph{From the Littlewood-Paley-Stein Inequality to
the Burkholder-Gundy Inequality}, arXiv:2111.05164; Trans. Amer. Math. Soc.
376 (2023), 371--389.
\item Q. Xu, \emph{Holomorphic functional calculus and vector-valued
Littlewood-Paley-Stein theory for semigroups}, arXiv:2105.12175.
\item Q. Xu, \emph{Optimal orders of the best constants in the
Littlewood-Paley inequalities}, arXiv:2105.12795; J. Funct. Anal. 283 (2022),
109570.
\item B. Davis, \emph{On the integrability of the martingale square function},
Israel J. Math. 8 (1970), 187--190.
\item D. L. Burkholder, \emph{Sharp inequalities for martingales and
stochastic integrals}, Asterisque 157--158 (1988), 75--94.
\end{enumerate}

\section*{Human Review Recommendation}

The packet is recommended for review as a partial positive endpoint theorem.
The main checks are: (i) confirm the indexing between \(a_n\), \(D_n\), and
\(\lambda_n\) in the Xu-Zhang semigroup; (ii) verify the Gram estimate in
Lemma~\ref{lem:gram}; and (iii) confirm the endpoint Davis inequality is being
used in the correct scalar martingale form.

\end{document}
